%% ---------------------------------------------------------------------
\section{Ablations and controls}\label{sec:ablations}

The benchmark's defensibility rests on knowing \emph{which} ingredient
carries the gain. We ablate every one.

\subsection{The target is the ingredient}
\begin{table}[t]
\caption{Target and mechanism controls on CIFAR-100 (champion
configuration except the varied ingredient; $\Delta$ vs.\ shared
baseline, three seeds). The gain needs the moment structure specifically:
random fixed targets do nothing, a $2\times$-cost learned teacher does
nothing anywhere on an eight-point envelope, and the classical HOG
descriptor recovers about half.}
\label{tab:controls}
\centering\scriptsize
\setlength{\tabcolsep}{2.6pt}
\resizebox{\columnwidth}{!}{%
\begin{tabular}{@{}lrrl@{}}
\toprule
aux target & @5\% & @10\% & note \\
\midrule
moment magnitude (ours) & $+3.31$ & $+2.81$ & phase-invariant energy \\
structure tensor &--& $+1.78$ & mild nonlinearity \\
steerable harmonics &--& $+1.01$ & principled rot.-inv. \\
rotation-invariant energy &--& $+0.84$ & \\
oriented edges (Gabor) &--& $-0.24$ & raw edges hurt \\
2nd-order invariants &--& $-2.97$ & over-processed \\
\midrule
random fixed maps & $+0.01$ & $+0.14$ & kills ``any aux signal'' \\
learned teacher (FitNets) & $-0.36$ & $+0.16$ & $2\times$ cost, ${\sim}0$ on 8-pt envelope \\
HOG target & $+1.22$ & $+1.37$ & descriptor helps; moments $2\times$ better \\
\midrule
cosine loss (not MSE) & -- & $+0.46$ & magnitude scale matters \\
\bottomrule
\end{tabular}}
\end{table}

Table~\ref{tab:controls}: a random fixed target of identical shape moves
accuracy by at most $0.14$ points, the mechanism is not ``regression as
regularization''. A FitNets teacher trained on the same data at twice the
cost does nothing anywhere we measured it (eight fractions, $-0.25$ to
$+0.73$), making the free hand-crafted target strictly better than an
expensive learned one in this regime. Among hand-crafted families,
phase-invariant \emph{magnitude} wins decisively; discarding magnitude
scale (cosine loss) forfeits most of the gain, and heavily processed
invariants are actively harmful targets.

\subsection{Placement, schedule, and head}
The tap depth is flat across the first three stages (within seed noise at
both 1\% and 10\%) with a cliff only at the final stage: the one design
knob in the study that is \emph{not} a data-regime knob. The schedule
start $\lambda_0$ \emph{is} a regime knob (best $2.0$ at 1--2\%, $1.0$ at
3--10\%, $0.3$ at 15--25\%), but its reach is bounded: three attempts to
push $\Delta$ past the measured feature gain by retuning $\lambda_0$ all
failed, consistent with the law, the knob cannot manufacture $G$ it can
only realize it. Head-norm projection, derived from a traced scale
degeneracy (the aux loss is invariant under feature-shrink with head-inflate,
which collapses bottleneck ResNets), converts ResNet-50 from a bistable
failure ($\{41.2, 36.4, 42.4\}$ across seeds) into a stable $+3.9$
($\{44.5, 45.0, 44.4\}$) and is free elsewhere; we adopt it always-on.
Notably, three plausible fixes tried \emph{before} tracing the mechanism
(gradient balancing, disabling AMP, BatchNorm on the tap) all failed or
backfired, the last making collapse $25\times$ faster, a small case
study in mechanism-first debugging.

\subsection{One configuration, every architecture and depth}
A single champion setting transplants without retuning: ResNet-18/34/50
land $+4.3/+4.0/+3.9$ at the same cell; the new depth sweep confirms one
$\lambda_0$ across ResNet-34/50 on six further domains ($+0.4$ to $+9.0$
at 5--10\% everywhere except CUB's deep-scarcity floor, decaying on the
right flank as the law requires). Head \emph{form} is irrelevant where it
plausibly might not have been: linear, cosine, and nearest-class-mean
readouts measure the same aux-vs-baseline gap to within $0.3$ points,
the left-flank bottleneck is label information, not classifier
expressivity.

\subsection{The forward-path alternative}
Placing the same bank \emph{in the forward path}, concatenated fixed
channels at the stem, the classical hybrid design, yields the study's
sharpest structural contrast. Forward-path stems are the low-data
accuracy leaders (energy-magnitude reaches $+2.6$ and $+3.5$ at CIFAR-100
1--2\%, beating the auxiliary prior there), but every fixed forward
variant we test, linear Gabor at two kernel scales, and four nonlinear
energy families, pays a penalty band at 10--25\% data ($-1.4$ to
$-7.3$) that no filter choice escapes, and the phase-invariant variants
never wash out even at 100\% ($-3.9$): a hard constraint permanently
occupies input bandwidth that abundant data would rather use. The
auxiliary placement of the \emph{same} knowledge is what removes the
band, which is the design argument for fusing priors through
training-only targets rather than through architecture.

\subsection{Bank width}
Widening the pinned 8-pair bank to 12 channels (via a third octave or two
extra orientations, indistinguishable at $10{\times}10$ seeds) buys a
real but local $+0.3$--$0.5$ at Tiny-ImageNet's 64\,px and nothing at
32\,px; across five further 64\,px domains the pooled excess fails our
pre-registered re-pin criterion. The committed 8-pair bank therefore
stays, with cross-domain evidence rather than convenience behind it.

\subsection{Stabilization as a second, distinct effect}
Across ConvNeXt (under the frozen SGD recipe), ResNet-50 (without
head-norm), Swin-T (on five datasets), and Swin at ImageNet64, baselines
are bistable, some seeds train, some sit at chance, while the prior
arm trains tightly every time. Swin's probe variance shows the
instability lives in the \emph{features} (baseline probe $\sigma$ up to
$4.3$ vs.\ $0.2$ with the prior). We report stabilization as a real,
separate property; we do not claim routine variance reduction (a
$10{\times}10$-seed test of that claim came back negative, and we report
that too).

%% ---------------------------------------------------------------------
\section{A decision guide}\label{sec:guide}

\begin{table}[t]
\caption{What to use when: the grid distilled to recommendations. Costs
are training-compute multiples of the baseline; gains are typical measured
ranges, not promises. ``Prior'' is the champion MomentAux configuration
verbatim.}
\label{tab:guide}
\centering\scriptsize
\setlength{\tabcolsep}{2.6pt}
\resizebox{\columnwidth}{!}{%
\begin{tabular}{@{}p{0.30\linewidth}p{0.22\linewidth}lp{0.27\linewidth}@{}}
\toprule
regime & recommendation & cost & measured basis \\
\midrule
Conv backbone, photo-like data, mid band (2--15\%), $2\times$ affordable & SimCLR init & $2\times$ & $+2..{+}9$ over baseline; beats prior by $+2..{+}5$ \\
Conv, photo-like, $2\times$ \emph{not} affordable & prior & $1.02\times$ & $+1.4..{+}6.7$; ${\sim}$SSL parity at extremes \\
Conv, extreme scarcity (${\leq}1$--$2\%$) & prior (or fwd stem) & $1.02\times$ & SSL margin ${\leq}1$ pt; forward stem best ${\leq}2\%$ \\
Satellite-like statistics, low data & prior & $1.02\times$ & beats SimCLR at 7 of 11 EuroSAT fractions \\
Any ViT trained the modern way (DeiT aug) & prior \emph{and} augmentation & $1.02\times$ & flip: beats SSL at every fraction; $+14..{+}25$ over recipe alone \\
ViT and Swin from scratch, any scale & prior & $1.02\times$ & $+3.2$ at 1.28M imgs; $+13/{+}26$ at 224\,px; stabilizes Swin \\
ImageNet-transferable task & transfer; \textbf{never} add prior &--& tax $-10..{-}24$, feature-side \\
Domain-shifted from ImageNet (e.g.\ stains) & prior from scratch competitive & $1.02\times$ & transfer advantage shrinks to $+2..{+}3$ \\
Data sufficient (${\geq}50\%$) & nothing (all converge) &--& every intervention within $\pm 0.8$ \\
\bottomrule
\end{tabular}}
\end{table}

Table~\ref{tab:guide} is the benchmark's practical distillate. Its
strongest single row is the modern-ViT one: for anyone training a small
or mid-sized ViT with the standard recipe, a fixed spectral auxiliary
target at $2\%$ overhead is, in every regime we measured, either the best
or tied-best intervention available, and it composes with the recipe
rather than competing with it.

%% ---------------------------------------------------------------------
\section{Discussion}\label{sec:discussion}

\paragraph{Why a law, and why this one.}
The grid's size is not its point; its point is that one low-dimensional
structure explains it. The decomposition in Eq.~\ref{eq:law} says that an
intervention's visible benefit is feature gain passed through a
readout channel whose behavior is set by how well the task is already
being solved. This explains, with one mechanism, observations that
otherwise require separate stories: why gains collapse at extreme
scarcity even when features improve (readout deeply negative, five
labels per class cannot express a better representation); why relabelling
the same pixels with coarser classes triples the visible gain (the
baseline crosses the readout zero, not because features changed; we
measure them unchanged); why every intervention converges at sufficiency
(both terms go to zero); and why probes read ${\sim}2$ points more gain
under coarse evaluation than fine on the same checkpoints (the measuring
stick itself is label-hungry).

\paragraph{The currency account as fusion theory in miniature.}
Our stack, substitute, or tax taxonomy is, we suggest, the general shape of
knowledge--data fusion outcomes: fusion adds value exactly when the
sources' contributions are non-overlapping \emph{in feature space}, and
harms when one source's dynamics (here, early shaping) overwrite what the
other supplied. The practical corollary is that fusion decisions can be
made by measurement rather than fashion: a cheap linear probe of what
each candidate source contributes predicts whether their combination is
worth anything, before any joint training is run.

\paragraph{What the attention result does and does not say.}
The persistent, scale-growing ViT deficit is a statement about
\emph{training small-data vision transformers from scratch}, where we
show a fixed spectral target supplies structure that neither 1.28M images
nor the DeiT recipe nor $2\times$-compute SSL supplies as cheaply. It is
not a claim about web-scale pre-training, where data eventually buys
everything (our own convolutional cells show exactly that limit). The
right reading is deployment-oriented: wherever from-scratch ViTs are
trained on $10^4$--$10^6$ images, medical imaging, remote sensing,
industrial inspection, a $2\%$-overhead structural prior is close to
free accuracy, and the smaller-model-plus-prior configuration can beat
the larger model outright.

\paragraph{On being wrong in public.}
Four of our major pre-registered predictions missed: SSL did not collapse
at 1\% on convolutional backbones (it won everywhere under plain
augmentation); augmentation did not substitute for the prior (it
amplified it); stronger contrastive views did not help SimCLR (they hurt
it); and an early ``gain tracks deficit'' law was falsified outright and
retracted. We report these because the misses carry as much information
as the hits, each redirected the program, and because a benchmark
that cannot be wrong is not measuring anything.

%% ---------------------------------------------------------------------
\section{Limitations and future directions}\label{sec:limits}

\paragraph{Known boundary of the law.} Ten resolvable cells violate the
sign law, concentrated on Food-101 and PathMNIST at ${\geq}20\%$ data,
all with the same signature: the prior still improves frozen-feature
quality while costing end-to-end accuracy. The current law has no term
for this ``better features, worse accuracy'' regime; modelling it,
plausibly an interaction between late-stage optimization and
sufficiency, is the clearest open problem the grid poses.

\paragraph{Scope of the instrument.} The frozen recipe is, in effect, a
ResNet recipe: architectures that cannot train under it (ConvNeXt, Swin
under SGD) enter only through diagnostic AdamW cells, and their
bistability, itself a finding, limits headline claims. PathMNIST's
shifted test split confounds its high-data cells for every method. Probes
lose interpretability at 100\% cells (probe-ceiling) and at ImageNet64
are budget-bound by construction. One prior family (Gabor and moment energy)
is tested as the knowledge source; the taxonomy predicts, but does not
yet demonstrate, that other structural priors will slot into the same
currency logic.

\paragraph{Future directions.} Three follow directly. (i) \emph{The
exception cluster}: instrument the late-training dynamics of high-data
Food-101 and PathMNIST cells to locate where good features fail to become
accuracy. (ii) \emph{Other currencies}: repeat the fusion matrix with
priors carrying genuinely different structure (depth, symmetry,
frequency-phase) to test whether currency-overlap measured by probes
predicts stacking in general. (iii) \emph{Attention at deployment
scale}: the model-scale trend ($+13$ rising to $+26$ from ViT-S to ViT-B at fixed
data) invites the obvious next point on the curve, and the
stabilization-plus-structure account suggests trying the prior as a
warmup for large-ViT pre-training proper.

%% ---------------------------------------------------------------------
\section{Conclusion}\label{sec:conclusion}

We built an instrument, one recipe, committed data, pinned filters,
pre-registered predictions, and used it to measure how a nearly free
structural prior fuses with the modern menu of data-efficiency
interventions across 14 datasets, 7 architectures and four orders of
magnitude in data. The measurements organize under one empirical law
separating feature gain from readout, sharp enough to predict unseen
backbones and to survive ImageNet scale within a point; and they support
a three-way fusion taxonomy, stack, substitute, tax, decided by
whether the fused sources carry the same information currency. The
practical residue is a decision guide whose strongest entry is also the
study's headline: for vision transformers trained the way practitioners
actually train them, a fixed spectral target costing $2\%$ of training
compute and nothing at inference is worth between $+13$ and $+26$
accuracy points, and no $2\times$-compute alternative we measured matches
it.
